\fichatitulo{49 · BENCHMARK DE GROUNDING}{Findings ACL · 2025}{GaRAGe}
{\small Sorodoc et al.. \textit{GaRAGe}. Findings ACL · 2025. \href{https://aclanthology.org/2025.findings-acl.875/}{Fonte consultada}.\par}

\campo{Problema e objetivo}
Como avaliar grounding relevante e recusa quando as fontes não permitem responder?

\campo{Principal contribuição · síntese interpretativa}
Disponibiliza benchmark com respostas longas, passagens de grounding anotadas e perguntas que exigem deflexão.

\campo{Método / natureza da evidência}
2.366 perguntas, mais de 35 mil passagens e anotação profissional de relevância, respostas e suficiência.

\campo{Resultado ou argumento principal}
Nos testes do benchmark, factualidade sensível à relevância chega a 60%, F1 de atribuição a 58,9% e taxa de verdadeiros positivos em deflexões a 31%.

\campo{Excertos-chave · seleção literal}
\begin{itemize}
\excerto{1}{models tend to over-summarise}{Resumo.}
\end{itemize}

\campo{Limitações e análise crítica}
A composição de fontes web e privadas e a avaliação com modelos comerciais limitam a equivalência com o acervo do Senado. Os números são resultados do benchmark, não taxas gerais.

\campo{Aproveitamento na dissertação · proposta}
Metodologia: incluir perguntas não respondíveis, suficiência documental e atribuição a passagens.

\registro{Fonte consultada nesta preparação: PDF publicado; consulta focal do resumo, construção das anotações, resultados e limitações. Chave: \texttt{sorodocEtAl2025garage}.}
